% ===== PRÉAMBULE CANONIQUE — à coller VERBATIM en tête de CHAQUE .tex =====
\documentclass[11pt,a4paper]{article}
\usepackage[utf8]{inputenc}\usepackage[T1]{fontenc}\usepackage{lmodern}
\usepackage[french]{babel}
\usepackage{amsmath,amssymb,mathtools}\usepackage{icomma}
\usepackage{siunitx}\sisetup{locale=FR}  % \qty{}{}, \si{}, \ang{}
\usepackage{xcolor}\usepackage{colortbl}
\usepackage{tikz}\usepackage{pgfplots}\pgfplotsset{compat=1.18}
\usepackage[siunitx,europeanresistors]{circuitikz} % circuits (élec.)
\usepackage[most]{tcolorbox}\usepackage{enumitem}\usepackage{array,booktabs}
\usepackage[a4paper,margin=2.2cm]{geometry}
\usetikzlibrary{arrows.meta,calc,angles,quotes,positioning,patterns,%
  backgrounds,shapes.geometric,decorations.pathmorphing,decorations.markings,intersections}
\definecolor{primary}{RGB}{222,49,99}\definecolor{primarydark}{RGB}{150,22,55}
\definecolor{defcol}{RGB}{0,114,178}\definecolor{methcol}{RGB}{52,92,148}
\definecolor{excol}{RGB}{0,135,90}\definecolor{attncol}{RGB}{200,40,55}
\tcbset{boxbase/.style={enhanced,breakable,boxrule=0.9pt,arc=2.5pt,
  left=3mm,right=3mm,top=2.3mm,bottom=2.3mm,fonttitle=\bfseries,coltitle=white}}
% --- boîtes TUTORIEL (noms EXACTS, ne pas renommer) ---
\newtcolorbox{cle}[1][]{boxbase,colback=defcol!6,colframe=defcol,
  title={\textsc{À retenir}\ifx\relax#1\relax\relax\else\textmd{ : \itshape #1}\fi}}
\newtcolorbox{methode}[1][]{boxbase,colback=methcol!7,colframe=methcol,sharp corners,
  title={\textsc{Méthode}\ifx\relax#1\relax\relax\else\textmd{ --- \itshape #1}\fi}}
\newtcolorbox{exemplebox}[1][]{boxbase,colback=excol!5,colframe=excol,
  title={\textsc{Exemple}\ifx\relax#1\relax\relax\else\textmd{ : \itshape #1}\fi}}
\newtcolorbox{piege}[1][]{boxbase,colback=attncol!6,colframe=attncol,
  title={\textsc{Piège}\ifx\relax#1\relax\relax\else\textmd{ : \itshape #1}\fi}}
\newtcolorbox{remarque}[1][]{boxbase,colback=black!4,colframe=black!45,
  title={\textsc{Remarque}\ifx\relax#1\relax\relax\else\textmd{ : \itshape #1}\fi}}
% --- boîtes EXERCICE / CORRIGÉ (noms EXACTS, auto-comptées) ---
\newtcolorbox[auto counter]{exo}[1][]{boxbase,colback=primary!4,colframe=primary,
  title={\textsc{Exercice}~\thetcbcounter\ifx\relax#1\relax\relax\else\textmd{ --- \itshape #1}\fi}}
\newtcolorbox[auto counter]{corrige}[1][]{boxbase,colback=excol!4,colframe=excol,
  title={\textsc{Corrigé}~\thetcbcounter\ifx\relax#1\relax\relax\else\textmd{ --- \itshape #1}\fi}}
\newcommand{\vv}[1]{\vec{#1}}\newcommand{\ds}{\displaystyle}
\newcommand{\R}{\mathbb{R}}\newcommand{\N}{\mathbb{N}}\newcommand{\Z}{\mathbb{Z}}
% =========================================================================
\begin{document}
% THEME_TITLE: Intensité du courant, charge et loi des nœuds
\usetikzlibrary{babel}

\begin{center}
{\LARGE\bfseries\color{primarydark} Thème 1 — Intensité, charge et loi des nœuds}\\[2pt]
{\color{primary}\itshape Compter les charges, mesurer un courant, répartir les intensités}
\end{center}

\begin{cle}[intensité d'un courant continu]
Le \textbf{courant} est un déplacement ordonné de charges. Son \textbf{intensité} $I$ mesure la
quantité de charge $Q$ qui traverse une section du conducteur \emph{par unité de temps} :
\[ \boxed{\;I=\dfrac{Q}{t}\;}\qquad \si{\ampere}=\si{\coulomb\per\second}. \]
La charge transportée par $n$ électrons vaut $\boxed{Q=n\,e}$, où $e=\qty{1.6e-19}{\coulomb}$ est la
\textbf{charge élémentaire}. On en tire le nombre d'électrons : $n=\dfrac{Q}{e}$.
\end{cle}

\begin{cle}[sens conventionnel et mesure de l'intensité]
Le \textbf{sens conventionnel} du courant va de la borne $+$ vers la borne $-$ \emph{à l'extérieur}
du générateur ; les \textbf{électrons} circulent en sens \emph{inverse}. On mesure $I$ avec un
\textbf{ampèremètre}, toujours branché \textbf{en série} : le courant doit le traverser.
\end{cle}

\begin{center}
\begin{circuitikz}[scale=0.92,transform shape]
  \draw (0,0) to[battery1,l_=$G$] (0,2.3)
        to[closing switch,l=$K$] (2.6,2.3)
        to[lamp,l=$L$] (2.6,0)
        to[ammeter](0,0);
  \draw[-{Stealth[length=2mm]},primary,line width=1pt] (0.9,2.42) -- (1.6,2.42)
     node[midway,above,font=\footnotesize,primary]{$I$};
\end{circuitikz}\\[1pt]
{\footnotesize L'ampèremètre (cercle $A$) est \textbf{en série} : il indique l'intensité $I$ du circuit.}
\end{center}

\begin{cle}[loi des nœuds (première loi de Kirchhoff)]
Un \textbf{nœud} est un point où se rejoignent plusieurs fils. Aucune charge ne s'y accumule : la
somme des intensités qui \emph{arrivent} égale la somme de celles qui en \emph{partent} :
\[ \boxed{\;\sum I_{\text{arrivant}}=\sum I_{\text{partant}}\;} \]
\end{cle}

\begin{center}
\begin{circuitikz}[scale=0.9,transform shape]
  \coordinate (N) at (0,0); \coordinate (M) at (4.4,0);
  \draw[thick] (N) -- (0,1.4); \draw (0,1.4) to[R,l=$R_1$] (4.4,1.4); \draw[thick] (4.4,1.4) -- (M);
  \draw (N) to[R,l=$R_2$] (M);
  \draw[thick] (N) -- (0,-1.4); \draw (0,-1.4) to[R,l=$R_3$] (4.4,-1.4); \draw[thick] (4.4,-1.4) -- (M);
  \draw[-{Stealth[length=2mm]},primary,line width=1pt] (-1.5,0) -- (-0.12,0)
     node[midway,above,font=\footnotesize,primary]{$I$};
  \fill (N) circle(2pt); \node[below left,font=\footnotesize] at (N){$N$};
  \draw[-{Stealth[length=1.6mm]},primary] (0.15,0.55) -- (0.55,1.15) node[right,font=\scriptsize,primary]{$I_1$};
  \draw[-{Stealth[length=1.6mm]},primary] (0.7,0) -- (1.3,0) node[above,font=\scriptsize,primary]{$I_2$};
  \draw[-{Stealth[length=1.6mm]},primary] (0.15,-0.55) -- (0.55,-1.15) node[right,font=\scriptsize,primary]{$I_3$};
\end{circuitikz}\\[1pt]
{\footnotesize Au nœud $N$ : $I=I_1+I_2+I_3$ (trois récepteurs \textbf{en parallèle}).}
\end{center}

\begin{methode}[compter les électrons qui traversent une section]
\begin{enumerate}[leftmargin=1.7em,topsep=2pt,itemsep=1pt]
  \item \textbf{Charge} : $Q=I\,t$ (penser à convertir $t$ en secondes).
  \item \textbf{Nombre d'électrons} : $n=\dfrac{Q}{e}=\dfrac{I\,t}{e}$ avec $e=\qty{1.6e-19}{\coulomb}$.
\end{enumerate}
\end{methode}

\begin{piege}[série $\neq$ parallèle : ne pas confondre les deux règles]
\textbf{En série} (aucun nœud entre les dipôles) : l'intensité est \textbf{la même partout},
$I=I_1=I_2=\dots$. \textbf{En parallèle} (présence de nœuds) : les intensités \textbf{s'ajoutent},
$I=I_1+I_2+\dots$. Si un fil est \emph{coupé} sur une branche, l'intensité y est \textbf{nulle}.
\end{piege}

\begin{remarque}[courant variable : intensité instantanée]
Si la charge ne traverse pas la section à débit constant, le courant est \emph{variable} et l'on
définit l'\textbf{intensité instantanée} $i(t)=\dfrac{\mathrm{d}q}{\mathrm{d}t}$ (dérivée de la
charge). C'est le cas du \emph{courant alternatif} du secteur, dont l'intensité change même de signe.
\end{remarque}

\end{document}
